% ==============================================================================
% Modelo para Anteprojeto de Graduação - Ciência da Computação (em Português)
% Prof. Vítor E. Silva Souza - NEMO / PPGI | DI / CT / UFES
% Com adaptações feitas pelo Colegiado de Engenharia de Computação / CT / UFES e pela profª. Monalessa Perini Barcellos (NEMO/UFES).
%
% Baseado em abtex2-modelo-trabalho-academico.tex, v-1.9.2 laurocesar
% Copyright 2012-2014 by abnTeX2 group at http://abntex2.googlecode.com/
%
% This work may be distributed and/or modified under the conditions of the LaTeX
% Project Public License, either version 1.3 of this license or (at your option)
% any later version. The latest version of this license is in
% http://www.latex-project.org/lppl.txt.
% ==============================================================================

% Usa o estilo abntex2, configurando detalhes de formatação e hifenização.
\documentclass[
	12pt,				% Tamanho da fonte.
	openright,			% Capítulos começam em página ímpar (insere página vazia caso preciso).
	oneside,			% Para impressão em verso e anverso. Oposto a oneside.
	a4paper,			% Tamanho do papel.
	english,			% Idioma adicional para hifenização.
	french,				% Idioma adicional para hifenização.
	spanish,			% Idioma adicional para hifenização.
	brazil				% O último idioma é o principal do documento.
	]{abntex2}



%%% Importação de pacotes. %%%

% Pacotes marcados com "%(desativado)" não são usados no texto e foram desligados para a
% compilação ficar mais rápida. Para usar um deles, basta tirar o "%" da linha.

% Conserta o erro "No room for a new \count" em versões antigas do LaTeX; nas atuais, já
% vem embutido.
%(desativado) \usepackage{etex}
%\reserveinserts{28}

% Usa a fonte Latin Modern.
\usepackage{lmodern}

% Seleção de códigos de fonte.
\usepackage[T1]{fontenc}

% Codificação do documento em Unicode.
\usepackage[utf8]{inputenc}

% Usado pela ficha catalográfica.
\usepackage{lastpage}

% Indenta o primeiro parágrafo de cada seção.
\usepackage{indentfirst}

% Controle das cores.
\usepackage[dvipsnames]{xcolor}

% Inclusão de gráficos.
\usepackage{graphicx}

% Tabularx package: melhor controle de leiaute de tabelas.
\usepackage{tabularx}

% Inclusão de páginas em PDF diretamente no documento (para uso nos apêndices).
%(desativado) \usepackage{pdfpages}

% Para melhorias de justificação. A expansão de fontes fica desligada: é a parte mais lenta do
% pacote e pesa no limite de tempo de compilação do Overleaf gratuito; a margem óptica continua.
\usepackage[expansion=false]{microtype}

% Citações padrão ABNT.
\usepackage[brazilian,hyperpageref]{backref}
\usepackage[alf]{abntex2cite}
\renewcommand{\backrefpagesname}{Citado na(s) página(s):~}		% Usado sem a opção hyperpageref de backref.
\renewcommand{\backref}{}										% Texto padrão antes do número das páginas.
\renewcommand*{\backrefalt}[4]{									% Define os textos da citação.
	\ifcase #1
		Nenhuma citação no texto.
	\or
		Citado na página #2.
	\else
		Citado #1 vezes nas páginas #2.
	\fi}

% \rm is deprecated and should not be used in a LaTeX2e document
% http://tex.stackexchange.com/questions/151897/always-textrm-never-rm-a-counterexample
\renewcommand{\rm}{\textrm}

% Inclusão de símbolos não padrão.
\usepackage{amssymb}
%(desativado) \usepackage{eurosym}

% Para utilizar \eqref para referenciar equações.
\usepackage{amsmath}

% Permite mostrar figuras muito largas em modo paisagem com \begin{sidewaysfigure} ao invés de \begin{figure}.
%(desativado) \usepackage{rotating}

% Permite customizar listas enumeradas/com marcadores.
\usepackage{enumitem}

% Notas de rodapé avançadas: sem uso no texto (as notas comuns não precisam dele).
%(desativado) \usepackage{bigfoot}

% Permite inserir hiperlinks com \url{}.
\usepackage{hyperref}

% Permite usar o comando \hl{} para evidenciar texto com fundo amarelo. Útil para chamar atenção a itens a fazer.
%(desativado) \usepackage{soulutf8}

% Colorinlistoftodos package: to insert colored comments so authors can collaborate on the content.
% É o pacote mais pesado da lista, porque carrega o TikZ; os três comandos andam juntos.
%(desativado) \usepackage[colorinlistoftodos, textwidth=20mm, textsize=footnotesize]{todonotes}
%(desativado) \newcommand{\aluno}[1]{\todo[author=\textbf{Aluno},color=green!30,caption={},inline]{#1}}
%(desativado) \newcommand{\professor}[1]{\todo[author=\textbf{Professor},color=red!30,caption={},inline]{#1}}

% Permite inserir espaço em branco condicional (incluído no texto final só se necessário) em macros.
%(desativado) \usepackage{xspace}

% Permite incluir listagens de código com o comando \lstinputlisting{}.
\usepackage{listings}
\usepackage{caption}
\DeclareCaptionFont{white}{\color{white}}
\DeclareCaptionFormat{listing}{\colorbox{gray}{\parbox{\textwidth}{#1#2#3}}}
\captionsetup[lstlisting]{format=listing,labelfont=white,textfont=white}
\renewcommand{\lstlistingname}{Listagem}
\definecolor{mygray}{rgb}{0.5,0.5,0.5}
\lstset{
	basicstyle=\scriptsize,
	breaklines=true,
	numbers=left,
	numbersep=5pt,
	numberstyle=\tiny\color{mygray},
	rulecolor=\color{black},
	showstringspaces=false,
	tabsize=2,
    inputencoding=utf8,
    extendedchars=true,
    literate=%
    {é}{{\'{e}}}1
    {è}{{\`{e}}}1
    {ê}{{\^{e}}}1
    {ë}{{\¨{e}}}1
    {É}{{\'{E}}}1
    {Ê}{{\^{E}}}1
    {û}{{\^{u}}}1
    {ù}{{\`{u}}}1
    {â}{{\^{a}}}1
    {à}{{\`{a}}}1
    {á}{{\'{a}}}1
    {ã}{{\~{a}}}1
    {Á}{{\'{A}}}1
    {Â}{{\^{A}}}1
    {Ã}{{\~{A}}}1
    {ç}{{\c{c}}}1
    {Ç}{{\c{C}}}1
    {õ}{{\~{o}}}1
    {ó}{{\'{o}}}1
    {ô}{{\^{o}}}1
    {Õ}{{\~{O}}}1
    {Ó}{{\'{O}}}1
    {Ô}{{\^{O}}}1
    {î}{{\^{i}}}1
    {Î}{{\^{I}}}1
    {í}{{\'{i}}}1
    {Í}{{\~{Í}}}1
}

% Posição exata das figuras com [H]: os marcadores de figura ficam onde estão no texto,
% sem se acumular na fila de figuras flutuantes.
\usepackage{float}

% Nota provisória embaixo de cada figura: o que a imagem deve mostrar. Sai antes da entrega.
% Para sumir com todas de uma vez, deixar o comando vazio: \newcommand{\notafigura}[1]{}
\newcommand{\notafigura}[1]{\par\smallskip{\footnotesize\itshape\color{red!75!black}[#1]\par}}

% Marcador das figuras que ainda não têm imagem: o brasão, guardado uma vez só (o PDF não repete o
% arquivo a cada figura). Ao entrar a imagem real, a linha vira um \includegraphics.
\newsavebox{\caixabrasao}
\AtBeginDocument{\sbox{\caixabrasao}{\includegraphics[width=0.16\textwidth]{figuras/brasao.jpg}}}
\newcommand{\marcadorbrasao}{\usebox{\caixabrasao}}

% Captura de tela com um filete em volta: o fundo claro do site se perderia no papel branco.
% Uso: \captura{figuras/arquivo.jpg} ou \captura[0.6\textwidth]{figuras/arquivo.jpg}.
\newcommand{\captura}[2][\textwidth]{{\setlength{\fboxsep}{0pt}\setlength{\fboxrule}{0.4pt}\fbox{\includegraphics[width=\dimexpr#1-0.8pt\relax]{#2}}}}



% Compilar só um capítulo, para caber no tempo do Overleaf enquanto se escreve: tire o "%" da
% linha abaixo e deixe nela só o capítulo em edição. Os outros mantêm a numeração da última
% compilação completa. Volte a comentar a linha para compilar o documento inteiro.
% \includeonly{capitulos/ch4-edugraph}



%%% Dados do trabalho. %%%

% TODO(autor): título provisório, o mesmo do anteprojeto. O título definitivo depende do
% enquadramento do trabalho e deve mencionar também a aplicação desenvolvida.
\titulo{Pensamento Computacional na Educação Básica: Um Mapeamento Sistemático sobre Estratégias, Impactos e Instrumentos de Avaliação para a Construção de um Repositório Educacional}
\autor{Lorenzo Rizzi Fiorot}
\local{Vitória, ES}
\data{2026}

% Rótulos no feminino: a orientação é de duas professoras. Os comandos vêm antes de
% \orientador e \coorientador, que guardam o rótulo no momento em que são chamados.
\renewcommand{\orientadorname}{Orientadora:}
\renewcommand{\coorientadorname}{Coorientadora:}
\orientador{Profa. Dra. Maria Claudia Silva Boeres}
\coorientador{Profa. Dra. Rosiane da Silva Sangali}

\instituicao{
  Universidade Federal do Espírito Santo -- UFES
  \par
  Centro Tecnológico
  \par
  Colegiado do Curso de Ciência da Computação}
\tipotrabalho{Monografia (PG)}

% Natureza do trabalho, na fórmula das monografias do curso.
\preambulo{Monografia apresentada ao Curso de Ciência da Computação do Centro Tecnológico da Universidade Federal do Espírito Santo, como requisito parcial para obtenção do Grau de Bacharel em Ciência da Computação.}




%%% Configurações finais de aparência. %%%

% Altera o aspecto da cor azul.
\definecolor{blue}{RGB}{41,5,195}

% Informações do PDF.
\makeatletter
\hypersetup{
	pdftitle={\@title},
	pdfauthor={\@author},
	pdfsubject={\imprimirpreambulo},
	pdfcreator={LaTeX with abnTeX2},
	colorlinks=true,				% Colore os links (ao invés de usar caixas).
	linkcolor=blue,					% Cor dos links.
	citecolor=blue,					% Cor dos links na bibliografia.
	filecolor=magenta,				% Cor dos links de arquivo.
	urlcolor=blue,					% Cor das URLs.
	bookmarksdepth=4
}
\makeatother

% Espaçamentos entre linhas e parágrafos.
\setlength{\parindent}{1.3cm}
\setlength{\parskip}{0.2cm}



%%% Páginas iniciais do documento: capa, folha de rosto, ficha, resumo, tabelas, etc. %%%

% Inicia o documento.
\begin{document}

% Brasão da instituição.
\begin{figure}
	\centering
	\includegraphics[width=.20\textwidth]{figuras/brasao.jpg}
	\label{fig-brasao}
\end{figure}

\begin{center}
	\textbf{\textsf{UNIVERSIDADE FEDERAL DO ESPÍRITO SANTO}}

	\textbf{\textsf{CENTRO TECNOLÓGICO}}

	\textbf{\textsf{COLEGIADO DO CURSO DE CIÊNCIA DA COMPUTAÇÃO}}

	\large{\textbf{\textsf{  }}}

	\large{\textbf{\textsf{  }}}
\end{center}

% Retira espaço extra obsoleto entre as frases.
\frenchspacing

% Capa do trabalho.
\imprimircapa

% Folha de rosto (o * indica que haverá a ficha bibliográfica).
\imprimirfolhaderosto*


% Ficha catalográfica, composta no próprio documento.
% Se a biblioteca fornecer a ficha em PDF, ela substitui a de baixo:
% \begin{fichacatalografica}
%     \includepdf{include-fichacatalografica.pdf}
% \end{fichacatalografica}
\begin{fichacatalografica}
	\vspace*{\fill}
	\hrule
	\begin{center}
	\begin{minipage}[c]{12.5cm}

	\imprimirautor

	\hspace{0.5cm} \imprimirtitulo  / \imprimirautor. --
	\imprimirlocal, \imprimirdata-

	\hspace{0.5cm} \pageref{LastPage} p. : il. (algumas color.) ; 30 cm.\\

	\hspace{0.5cm} \imprimirorientadorRotulo~\imprimirorientador\\

	\hspace{0.5cm} \imprimircoorientadorRotulo~\imprimircoorientador\\

	\hspace{0.5cm}
	\parbox[t]{\textwidth}{\imprimirtipotrabalho~--~\imprimirinstituicao,
	\imprimirdata.}\\

	\hspace{0.5cm}
		% TODO(autor): palavras-chave, as mesmas do resumo.
		1. Texto.
		I. Boeres, Maria Claudia Silva.
		II. Sangali, Rosiane da Silva.
		III. Universidade Federal do Espírito Santo.
		IV. \imprimirtitulo \\

	% TODO(autor): o código CDU abaixo é o do modelo; a biblioteca informa o correto.
	\hspace{8.75cm} CDU 02:141:005.7\\

	\end{minipage}
	\end{center}
	\hrule
\end{fichacatalografica}


% Folha de aprovação.
% Depois da defesa, a cópia digitalizada da folha assinada pode substituir a de baixo:
% \includepdf{include-folhadeaprovacao.pdf}
\begin{folhadeaprovacao}
  \begin{center}
    {\ABNTEXchapterfont\large\imprimirautor}
    \vspace*{\fill}\vspace*{\fill}
    \begin{center}
      \ABNTEXchapterfont\bfseries\Large\imprimirtitulo
    \end{center}
    \vspace*{\fill}
    \hspace{.45\textwidth}
    \begin{minipage}{.5\textwidth}
        \imprimirpreambulo
    \end{minipage}%
    \vspace*{\fill}
   \end{center}
   % TODO(autor): data da defesa e membros da banca, com a instituição de cada um.
   Trabalho aprovado. \imprimirlocal, Texto.:
   \assinatura{\textbf{\imprimirorientador} \\ Orientadora}
   \assinatura{\textbf{\imprimircoorientador} \\ Coorientadora}
   \assinatura{\textbf{Texto.} \\ Texto.}
   \assinatura{\textbf{Texto.} \\ Texto.}
   \begin{center}
    \vspace*{0.5cm}
    {\large\imprimirlocal}
    \par
    {\large\imprimirdata}
    \vspace*{1cm}
  \end{center}
\end{folhadeaprovacao}


% Dedicatória (opcional).
% TODO(autor): escrever a dedicatória ou retirar este elemento.
\begin{dedicatoria}
   \vspace*{\fill}
   \centering
   \noindent
   \textit{Texto.} \vspace*{\fill}
\end{dedicatoria}


% Agradecimentos (opcional).
% TODO(autor): escrever os agradecimentos ou retirar este elemento.
\begin{agradecimentos}
Texto.
\end{agradecimentos}


% Epígrafe (opcional).
% TODO(autor): escrever a epígrafe ou retirar este elemento.
\begin{epigrafe}
    \vspace*{\fill}
	\begin{flushright}
		\textit{Texto.}
	\end{flushright}
\end{epigrafe}


% Resumo em português.
% TODO(autor): o resumo e as palavras-chave são escritos por último, com o texto pronto.
\setlength{\absparsep}{18pt}
\begin{resumo}
Texto.

\textbf{Palavras-chave}: Texto.
\end{resumo}

% Insere lista de ilustrações.
\pdfbookmark[0]{\listfigurename}{lof}
\listoffigures*
\cleardoublepage

% Insere lista de tabelas.
\pdfbookmark[0]{\listtablename}{lot}
\listoftables*
\cleardoublepage

% Lista de abreviaturas e siglas: só as que aparecem no texto, em ordem alfabética.
\begin{siglas}
  % As siglas comentadas eram usadas só no mapeamento (capítulo 3), que voltou a ser
  % esqueleto; voltam junto com o texto dele.
  % As de modelos de grafo (IRI, RDF, SPARQL, URI) saíram com a antiga seção de grafos de
  % conhecimento; a do W3C volta com a seção de acessibilidade.
  \item[API] \textit{Application Programming Interface}
  \item[AVA] Ambiente Virtual de Aprendizagem
  \item[BNCC] Base Nacional Comum Curricular
  %\item[CE] Critério de Exclusão
  \item[CEB] Câmara de Educação Básica
  %\item[CI] Critério de Inclusão
  \item[CIF] Comissão Intergovernamental de Financiamento para a Educação Básica de Qualidade
  \item[CNE] Conselho Nacional de Educação
  \item[CSS] \textit{Cascading Style Sheets}
  \item[Fundeb] Fundo de Manutenção e Desenvolvimento da Educação Básica e de Valorização dos Profissionais da Educação
  \item[HTML] \textit{HyperText Markup Language}
  \item[HTTP] \textit{Hypertext Transfer Protocol}
  %\item[IRI] \textit{Internationalized Resource Identifier}
  \item[JSON] \textit{JavaScript Object Notation}
  \item[MEC] Ministério da Educação
  \item[MSL] Mapeamento Sistemático da Literatura
  %\item[PAPC] Protocolo para Avaliação por Critérios do Pensamento Computacional
  \item[PC] Pensamento Computacional
  %\item[PICOC] \textit{Population, Intervention, Comparison, Outcome, Context}
  %\item[QP] Questão de Pesquisa
  %\item[RDF] \textit{Resource Description Framework}
  %\item[RENOTE] Revista Novas Tecnologias na Educação
  %\item[SBC] Sociedade Brasileira de Computação
  %\item[SOL] SBC OpenLib
  %\item[SOLO] \textit{Structure of the Observed Learning Outcome}
  %\item[SPARQL] \textit{SPARQL Protocol and RDF Query Language}
  %\item[STEAM] \textit{Science, Technology, Engineering, Arts and Mathematics}
  %\item[URI] \textit{Uniform Resource Identifier}
  \item[VAAR] Valor anual por aluno decorrente da complementação-VAAR do Fundeb
  \item[WCAG] \textit{Web Content Accessibility Guidelines}
  %\item[W3C] \textit{World Wide Web Consortium}
\end{siglas}

% Insere o sumário.
\pdfbookmark[0]{\contentsname}{toc}
\tableofcontents*
\cleardoublepage



%%% Início da parte de conteúdo do documento. %%%

% Marca o início dos elementos textuais.
\textual

% Capítulos, um arquivo por capítulo na pasta capitulos/.
% ==============================================================================
% PG - Lorenzo Rizzi Fiorot
% Capítulo 1 - Introdução
% ==============================================================================
\chapter{Introdução}
\label{sec-intro}

Com a incorporação da Computação à Base Nacional Comum Curricular (BNCC), o desenvolvimento do Pensamento Computacional (PC) passou a integrar formalmente o currículo da educação básica brasileira, previsto para ocorrer de forma transversal às demais áreas do conhecimento \cite{brasil2022computacao, brasil2022parecer}. Essa transversalidade atribui ao professor a tarefa de articular os conteúdos de seu próprio componente curricular aos conceitos computacionais; tais conexões, contudo, permanecem implícitas nos documentos normativos e dispersas na literatura científica, sem um meio didático que as torne acessíveis ao docente.

% TODO(autor): parágrafo de apresentação do trabalho. O que o repositório relaciona
% depende de como o trabalho enquadra o Pensamento Computacional: hoje a base de dados
% traz a estrutura curricular da BNCC (ano, área, componente, unidade temática, objeto de
% conhecimento e habilidade), sem conceitos do Pensamento Computacional. A aplicação
% atual é uma interface em React com uma API própria em Python sobre o banco Neo4j.
Texto.


%%% Início de seção. %%%
\section{Motivação e Justificativa}
\label{sec-intro-motjus}

A crescente presença das tecnologias digitais na vida cotidiana consolidou o entendimento de que competências relacionadas à Computação tornaram-se tão fundamentais quanto as aprendizagens tradicionalmente associadas à escolarização, como ler, escrever e calcular. No Brasil, a formalização desse entendimento ocorreu com as Normas sobre Computação na Educação Básica, que constituem complemento à BNCC e foram instituídas pela Resolução CNE/CEB nº 1, de 4 de outubro de 2022, da Câmara de Educação Básica (CEB) do Conselho Nacional de Educação (CNE) \cite{brasil2022computacao}. Esse documento normativo estabelece diretrizes para a inserção da Computação nos currículos da Educação Básica, orientando o desenvolvimento de conhecimentos, habilidades e competências relacionados às tecnologias digitais e determinando sua implementação pelos sistemas e pelas instituições de ensino brasileiras. Sua organização curricular estrutura-se em três eixos articulados: o \textbf{Pensamento Computacional}, o \textbf{Mundo Digital} e a \textbf{Cultura Digital} \cite{brasil2022complemento}. Neste trabalho, esse conjunto normativo é referido, por concisão, como BNCC Computação.

O modelo pedagógico adotado pela BNCC Computação não restringe a Computação a um componente curricular isolado. Ao contrário, prevê seu desenvolvimento de forma transversal e integrada às demais áreas do conhecimento. Essa articulação pressupõe que conteúdos já presentes em componentes como Matemática, Ciências e Língua Portuguesa possam servir de contexto para a introdução de conceitos computacionais.

A implementação, prevista para começar em até um ano após a homologação de 2022 \cite{brasil2022computacao} e reforçada por normas posteriores de integração curricular da educação digital \cite{brasil2025cne}, conta ainda com um mecanismo de incentivo financeiro: o alinhamento dos referenciais curriculares à BNCC, que inclui a Computação, integra as condicionalidades da complementação-VAAR do Fundo de Manutenção e Desenvolvimento da Educação Básica e de Valorização dos Profissionais da Educação (Fundeb), complementação de que decorre o valor anual por aluno (VAAR) \cite{brasil2020fundeb}. A Resolução CIF nº 15, de 12 de junho de 2025, da Comissão Intergovernamental de Financiamento para a Educação Básica de Qualidade (CIF), restringe a um único exercício a dispensa dessa exigência:

\begin{citacao}
Caso os referenciais curriculares não contemplem a Computação na Educação Básica --- Complemento à BNCC, a rede de ensino não será inabilitada em 2025 para fins de recebimento dos recursos da complementação do VAAR em 2026, devendo providenciar a adequação, de forma que tal situação não implique a inabilitação nos anos subsequentes \cite[art.~4º, §~3º]{brasil2025cif}.
\end{citacao}

\noindent Encerrada essa transição, as redes que não incorporarem a Computação a seus currículos ficam sujeitas à inabilitação e à perda na complementação-VAAR nos exercícios seguintes. Para além da dimensão financeira, a não implementação priva os estudantes de competências que a própria norma reconhece como essenciais ao exercício pleno da cidadania \cite{brasil2022parecer}.

Apesar da obrigatoriedade, a implementação esbarra em uma dificuldade concreta no âmbito da prática docente. Grande parte dos professores da educação básica não teve, em sua formação, contato com programação ou com o PC, o que faz da formação docente uma condição necessária --- e frequentemente ausente --- para a efetivação dessas competências \cite{boeres2024}. Mais especificamente, o modelo transversal exige que o professor reconheça de que maneira o seu próprio componente curricular pode servir de veículo para o ensino do PC; contudo, as conexões entre as habilidades previstas na BNCC e os conceitos computacionais permanecem implícitas e dispersas na documentação oficial, sem um meio didático e acessível que torne evidente o que, dentro de cada disciplina, pode ser aproveitado para ensinar PC de forma integrada --- o que transfere ao docente o ônus de estabelecer, por conta própria, articulações que demandam conhecimento especializado.

A literatura nacional já documenta diversas iniciativas de ensino do PC na educação básica, que vão do emprego da Computação Desplugada \cite{silvamadeira2024}, isto é, de atividades que dispensam o uso de computadores \cite[p.~50]{brackmann2017}, ao uso de ferramentas e ambientes de programação alinhados à BNCC \cite{osorio2025}, além de mapeamentos sistemáticos que caracterizam o campo e seus instrumentos de avaliação \cite{rege2023}. Esses trabalhos, contudo, concentram-se majoritariamente em relatos de intervenções pontuais --- associadas a ferramentas, séries ou contextos específicos --- ou em sínteses do estado da arte, sem disponibilizar um artefato reutilizável que torne explícitas, de forma navegável, as conexões entre os componentes curriculares da BNCC e os conceitos do PC. Soma-se a isso a escassez e a dispersão de materiais de apoio ao professor voltados a essa articulação, o que evidencia a necessidade de dar continuidade à investigação do tema sob uma perspectiva que não apenas caracterize as práticas existentes, mas também organize e disponibilize esse conhecimento de maneira acessível ao docente.

% TODO(autor): parágrafo que apresenta a contribuição diante da lacuna. No anteprojeto, ele
% prometia tornar navegáveis as relações entre os componentes curriculares da BNCC e os
% pilares do Pensamento Computacional; a redação depende do mesmo enquadramento do
% parágrafo de apresentação, no início do capítulo.
Texto.


%%% Início de seção. %%%
\section{Objetivos}
\label{sec-intro-obj}

% TODO(autor): o objetivo geral e os específicos do anteprojeto falavam em relacionar a
% BNCC ao Pensamento Computacional, em integrar a interface a um endpoint SPARQL e em
% especificar requisitos e avaliar a usabilidade. Reescrever conforme o enquadramento do
% trabalho e o que de fato foi entregue.
\subsection{Objetivo Geral}
\label{sec-intro-obj-geral}

Texto.

\subsection{Objetivos Específicos}
\label{sec-intro-obj-especificos}

Texto.


%%% Início de seção. %%%
\section{Escopo e Divisão de Responsabilidades}
\label{sec-intro-escopo}

% TODO(autor): no anteprojeto, a modelagem do grafo e o servidor de consultas cabiam à
% orientação, e o armazenamento e a consulta ficavam fora do escopo. Hoje a modelagem do
% domínio segue com a orientação, mas este trabalho desenvolveu também a API e mantém a
% base de dados, executando comandos Cypher validados. A redação desta divisão é do autor.
Texto.


%%% Início de seção. %%%
\section{Método de Desenvolvimento do Trabalho}
\label{sec-intro-met}

% TODO(autor): o método do anteprojeto (mapeamento, requisitos, prototipação,
% implementação e integração, testes) é reescrito no passado quando o escopo e os
% objetivos estiverem definidos.
Texto.


%%% Início de seção. %%%
\section{Organização da Monografia}
\label{sec-intro-organizacao}

Texto.

% ==============================================================================
% PG - Lorenzo Rizzi Fiorot
% Capítulo 2 - Referencial Teórico e Tecnologias Utilizadas
% ==============================================================================
\chapter{Referencial Teórico e Tecnologias Utilizadas}
\label{sec-referencial}

% TODO(autor): no anteprojeto, a abertura do capítulo também apresentava o repositório, como
% relação entre a BNCC e os pilares do Pensamento Computacional; essa apresentação depende
% do enquadramento do trabalho.
Este capítulo reúne os fundamentos teóricos do trabalho e as tecnologias empregadas em seu desenvolvimento. A Seção~\ref{sec-referencial-bncc} apresenta a estrutura da BNCC e da BNCC Computação, e a Seção~\ref{sec-referencial-pc}, o Pensamento Computacional. A Seção~\ref{sec-referencial-msl} caracteriza o mapeamento sistemático da literatura como método de pesquisa. A Seção~\ref{sec-referencial-delimitacao} delimita o conceito de repositório de conhecimento adotado. A Seção~\ref{sec-referencial-visualizacao} trata da visualização de grafos, e a Seção~\ref{sec-referencial-acessibilidade}, da acessibilidade na web. Por fim, a Seção~\ref{sec-referencial-tecnologias} descreve as tecnologias utilizadas.


%%% Início de seção. %%%
\section{A BNCC e a BNCC Computação}
\label{sec-referencial-bncc}

A BNCC é um ``documento de caráter normativo que define o conjunto orgânico e progressivo de aprendizagens essenciais que todos os alunos devem desenvolver ao longo das etapas e modalidades da Educação Básica'' \cite[p.~7]{brasil2018bncc}. O documento abrange a Educação Infantil, o Ensino Fundamental e o Ensino Médio e constitui referência nacional para a formulação dos currículos dos sistemas e das redes de ensino e das propostas pedagógicas das escolas \cite[p.~8--9]{brasil2018bncc}. As aprendizagens essenciais devem concorrer para o desenvolvimento de dez competências gerais, entendida a competência como ``a mobilização de conhecimentos (conceitos e procedimentos), habilidades (práticas, cognitivas e socioemocionais), atitudes e valores para resolver demandas complexas da vida cotidiana, do pleno exercício da cidadania e do mundo do trabalho'' \cite[p.~8]{brasil2018bncc}.

No Ensino Fundamental, etapa a que pertencem os dados do grafo apresentado neste trabalho, a BNCC organiza-se em cinco áreas do conhecimento --- Linguagens, Matemática, Ciências da Natureza, Ciências Humanas e Ensino Religioso ---, que reúnem os componentes curriculares \cite[p.~27]{brasil2018bncc}. Cada componente apresenta um conjunto de habilidades, que expressam as aprendizagens essenciais a serem asseguradas aos alunos. As habilidades relacionam-se a objetos de conhecimento, isto é, a conteúdos, conceitos e processos, e estes, por sua vez, organizam-se em unidades temáticas \cite[p.~28--29]{brasil2018bncc}. Uma unidade temática contempla uma gama maior ou menor de objetos de conhecimento, e cada objeto relaciona-se a um número variável de habilidades \cite[p.~29]{brasil2018bncc}. Esse encadeamento --- área do conhecimento, componente curricular, unidade temática, objeto de conhecimento e habilidade --- é a estrutura reproduzida no modelo do grafo.

Cada habilidade é identificada por um código alfanumérico \cite[p.~30]{brasil2018bncc}, cuja composição é apresentada na Tabela~\ref{tab:codigo-habilidade}, tomando-se como exemplo o código \texttt{EF01CI05}. Duas ressalvas do documento são relevantes para este trabalho: a numeração sequencial não representa uma ordem ou hierarquia esperada das aprendizagens, e a organização das habilidades em objetos de conhecimento e unidades temáticas expressa ``um arranjo possível (dentre outros)'', que não deve ser tomado como modelo obrigatório para o desenho dos currículos \cite[p.~31]{brasil2018bncc}. Ambas se estendem ao grafo, que reproduz esse arranjo.

\begin{table}[h]
    \centering
    \caption{Composição do código das habilidades do Ensino Fundamental}
    \label{tab:codigo-habilidade}
    \renewcommand{\arraystretch}{1.3}
    \begin{tabular}{|p{4.2cm}|p{6cm}|p{4.2cm}|}
        \hline
        \textbf{Parte do código} & \textbf{O que indica} & \textbf{No código \texttt{EF01CI05}} \\
        \hline
        Primeiro par de letras & Etapa de ensino & \texttt{EF}: Ensino Fundamental \\
        \hline
        Primeiro par de números & Ano ou bloco de anos & \texttt{01}: 1º ano \\
        \hline
        Segundo par de letras & Componente curricular & \texttt{CI}: Ciências \\
        \hline
        Último par de números & Posição na numeração sequencial do ano ou do bloco de anos & \texttt{05}: quinta habilidade \\
        \hline
    \end{tabular}
\end{table}

A Computação já era tratada na BNCC de 2018, que tematiza três dimensões da computação e das tecnologias digitais \cite[p.~473--474]{brasil2018bncc}: o pensamento computacional, discutido na Seção~\ref{sec-referencial-pc}; o mundo digital, que envolve as formas de processar, transmitir e distribuir a informação de maneira segura e confiável; e a cultura digital, voltada a uma participação mais consciente e democrática por meio das tecnologias digitais. No Ensino Fundamental, contudo, essas dimensões foram contempladas nas competências e habilidades dos diferentes componentes \cite[p.~474]{brasil2018bncc}, sem habilidades próprias, uma vez que a Computação não figura entre os componentes curriculares da etapa \cite[p.~27]{brasil2018bncc}.

O conjunto próprio de habilidades veio com a Resolução CNE/CEB nº 1, de 4 de outubro de 2022, que define normas sobre Computação na Educação Básica, em complemento à BNCC (art.~1º). A Resolução determina que o desenvolvimento e a formulação dos currículos considerem as tabelas de competências e habilidades anexas (art.~1º, §~2º), atribui aos Estados, aos Municípios e ao Distrito Federal a definição dos parâmetros e das abordagens pedagógicas de implementação (art.~2º) e prevê que o Ministério da Educação (MEC), em conjunto com esses entes, defina política para a formação docente, o apoio ao desenvolvimento de currículos e o apoio ao desenvolvimento de recursos didáticos compatíveis com as tabelas (art.~4º) \cite{brasil2022computacao}.

As tabelas constam do documento anexo ao Parecer CNE/CEB nº 2/2022 \cite{brasil2022complemento} e cobrem as três etapas da Educação Básica. No Ensino Fundamental, há sete competências e um conjunto de habilidades para cada ano, além de habilidades por etapa (do 1º ao 5º e do 6º ao 9º ano). Cada habilidade vincula-se a um objeto de conhecimento e a um dos três eixos --- Pensamento Computacional, Mundo Digital e Cultura Digital --- e vem acompanhada de explicação e de exemplos. Os códigos seguem a composição da Tabela~\ref{tab:codigo-habilidade}, com o par \texttt{CO} na posição do componente curricular: \texttt{EF01CO01}, por exemplo, identifica a primeira habilidade de Computação do 1º ano, ``Organizar objetos físicos ou digitais considerando diferentes características para esta organização, explicitando semelhanças (padrões) e diferenças'' \cite[p.~12]{brasil2022complemento}. A organização das tabelas é, portanto, análoga à dos demais componentes, com o eixo no lugar da unidade temática.

% TODO(autor): se o enquadramento do trabalho incluir a BNCC Computação no modelo do grafo,
% acrescentar aqui a relação entre as tabelas do complemento e o modelo.


%%% Início de seção. %%%
\section{Pensamento Computacional}
\label{sec-referencial-pc}

A popularização do termo Pensamento Computacional é atribuída ao artigo de \citeonline{wing2006}, como registram \citeonline[p.~562]{boeres2024}. Para a autora, trata-se de uma habilidade fundamental para todos, e não apenas para cientistas da computação, a ser acrescentada à leitura, à escrita e à aritmética na formação de toda criança; envolve a resolução de problemas, o projeto de sistemas e a compreensão do comportamento humano com base nos conceitos fundamentais da ciência da computação \cite[p.~33]{wing2006}. Nessa acepção, pensar computacionalmente significa conceituar, e não programar, e constitui um modo pelo qual as pessoas, e não os computadores, resolvem problemas \cite[p.~35]{wing2006}.

No contexto brasileiro, \citeonline{brackmann2017} propõe, a partir da fusão de diversas fontes, a seguinte definição:

\begin{citacao}
O Pensamento Computacional é uma distinta capacidade criativa, crítica e estratégica humana de saber utilizar os fundamentos da Computação, nas mais diversas áreas do conhecimento, com a finalidade de identificar e resolver problemas, de maneira individual ou colaborativa, através de passos claros, de tal forma que uma pessoa ou uma máquina possam executá-los eficazmente \cite[p.~29]{brackmann2017}.
\end{citacao}

\noindent A BNCC, por sua vez, afirma que o pensamento computacional ``envolve as capacidades de compreender, analisar, definir, modelar, resolver, comparar e automatizar problemas e suas soluções, de forma metódica e sistemática, por meio do desenvolvimento de algoritmos'' \cite[p.~474]{brasil2018bncc}. As três formulações tratam o PC como um modo de resolver problemas fundamentado na Computação, e as de \citeonline{wing2006} e \citeonline{brackmann2017} ressaltam que os processos envolvidos podem ser executados por uma pessoa ou por uma máquina. Por essa razão, o PC pode ser desenvolvido também por meio da Computação Desplugada, que dispensa o uso de computadores e recorre, por exemplo, a movimentos, cartões, desenhos e enigmas \cite[p.~50]{brackmann2017}.

Segundo \citeonline[p.~33]{brackmann2017}, o PC apoia-se em quatro pilares, interdependentes e voltados à resolução de problemas:

\begin{itemize}
    \item \textbf{decomposição}: identificar um problema complexo e dividi-lo em partes menores e mais fáceis de gerenciar;
    \item \textbf{reconhecimento de padrões}: analisar cada uma dessas partes em busca de problemas parecidos, já solucionados anteriormente;
    \item \textbf{abstração}: concentrar-se apenas nos detalhes importantes, ignorando as informações irrelevantes;
    \item \textbf{algoritmos}: criar passos ou regras simples para resolver cada um dos subproblemas encontrados.
\end{itemize}

\noindent Os mesmos pilares são adotados por \citeonline[p.~565]{boeres2024}, na análise da formação docente, e por \citeonline[p.~1752]{osorio2025}, no mapeamento entre o PC e as habilidades da BNCC.

Nos documentos curriculares, o PC aparece de duas formas. Na BNCC de 2018, as menções no Ensino Fundamental concentram-se na área de Matemática: os processos de resolução de problemas, de investigação, de desenvolvimento de projetos e de modelagem são considerados potencialmente ricos para o seu desenvolvimento \cite[p.~266]{brasil2018bncc}, e ressaltam-se a importância dos algoritmos e de seus fluxogramas e a identificação de padrões para estabelecer generalizações \cite[p.~271]{brasil2018bncc}. Na BNCC Computação, o PC constitui um eixo próprio, e os pilares manifestam-se nas habilidades desde os anos iniciais: no 1º ano, a identificação de semelhanças, ou padrões, na organização de objetos (\texttt{EF01CO01}) e a criação de sequências de passos, associadas à palavra ``Algoritmos'' (\texttt{EF01CO03}); no 2º ano, a identificação de atributos essenciais, própria da abstração, na modelagem de objetos (\texttt{EF02CO01}); no 3º ano, a decomposição de problemas complexos (\texttt{EF03CO03}) \cite[p.~12, 16, 18]{brasil2022complemento}.

% TODO(autor): parágrafo de fechamento, que liga o Pensamento Computacional ao repositório
% de conhecimento; depende do enquadramento do trabalho.


%%% Início de seção. %%%
\section{Mapeamento Sistemático da Literatura}
\label{sec-referencial-msl}

O Mapeamento Sistemático da Literatura (MSL) é uma metodologia de pesquisa secundária que visa identificar, categorizar e sintetizar o conjunto de evidências disponíveis sobre um determinado tema, distinguindo-se da revisão sistemática convencional por priorizar uma visão panorâmica da área em vez de responder a uma questão clínica ou empírica específica \cite{petersen2008systematic}. Enquanto a revisão sistemática tende a focar em perguntas de pesquisa bem delimitadas e a realizar sínteses quantitativas dos resultados, o mapeamento sistemático permite classificar e estruturar o estado da arte de forma mais abrangente, evidenciando lacunas, tendências e distribuição de estudos ao longo do tempo \cite{kitchenham2007guidelines, petersen2008systematic}. No presente trabalho, essa abordagem foi adotada para levantar o estado da arte sobre o ensino do Pensamento Computacional no contexto da educação básica brasileira, com ênfase na integração com os componentes curriculares da BNCC, fornecendo a base empírica necessária para a construção do repositório de conhecimento proposto. O protocolo, a condução e os resultados do mapeamento são apresentados no Capítulo~\ref{sec-msl}.


%%% Início de seção. %%%
\section{Delimitação Conceitual: Repositório de Conhecimento}
\label{sec-referencial-delimitacao}

A literatura técnica e científica emprega diferentes designações para soluções baseadas na web, cujas distinções conceituais são relevantes para a correta caracterização do artefato desenvolvido neste trabalho. Por \textit{repositório} entende-se, no campo da Ciência da Informação e da Arquivística, uma solução voltada ao armazenamento, ao gerenciamento e à preservação de longo prazo de materiais digitais, constituída por hardware, software, metadados e procedimentos normativos \cite{conarq2015}. Na Engenharia de Software, o termo designa ainda o local em que um sistema de controle de versão registra e mantém o histórico de alterações do código-fonte, como nos repositórios Git \cite{chacon2014}. Em ambas as acepções, a ênfase recai sobre a guarda e a organização de um acervo --- de conteúdo ou de código ---, e não sobre a sua mera apresentação.

Distinguem-se desse conceito outras designações correntes. Um \textit{site} (ou \textit{website}) é um conjunto de páginas web interligadas e hospedadas em um mesmo endereço de rede, acessadas a partir de uma página inicial \cite{w3c1999}, noção centrada na disponibilização de conteúdo e com interatividade tipicamente limitada. Um \textit{portal} designa um ponto único de acesso que agrega, em uma mesma interface, informações, recursos e serviços de fontes diversas, funcionando como porta de entrada para a utilização de serviços na internet \cite{gouveia2007}; sua característica definidora é a agregação, e não a estruturação de um acervo de conhecimento. Já uma \textit{plataforma}, em sua acepção mais ampla, corresponde a uma infraestrutura digital que possibilita a interação entre dois ou mais grupos de usuários \cite{srnicek2017}, conceito associado a ecossistemas extensíveis e de natureza sociotécnica, cuja aplicação exige rigor de definição \cite{dereuver2018}. O termo \textit{ambiente}, por sua vez, remete a um espaço de execução ou de interação e, no domínio educacional, consolidou-se na expressão \textit{ambiente virtual de aprendizagem} (AVA), que designa espaços digitais de mediação do processo de ensino-aprendizagem apoiados em tecnologias de informação e comunicação \cite{almeida2003}.

Cabe, por fim, distinguir os termos \textit{sistema web} e \textit{aplicação web} (\textit{webapp}). O primeiro designa um software hospedado em servidor e acessado por navegador, geralmente restrito a usuários cadastrados, em contraste com o site meramente informativo \cite{sommerville2011}. O segundo, mais preciso para soluções que integram interface de usuário (\textit{front-end}) e lógica de servidor (\textit{back-end}), compreende desde sites completos até funcionalidades especializadas e aplicações de processamento de dados \cite{pressman2011}.

Diante dessas distinções, este trabalho adota o termo \textit{repositório de conhecimento} para designar o artefato central produzido, uma vez que nenhuma das demais designações (\textit{site}, \textit{portal}, \textit{plataforma} ou \textit{ambiente}) captura adequadamente a natureza de acervo organizado e semanticamente explícito da solução. Esse repositório materializa-se na forma de um grafo de conhecimento. O termo \textit{aplicação web} é empregado, de modo complementar e não substitutivo, para designar especificamente a camada de software que disponibiliza o repositório, reservando-se o termo \textit{repositório} ao acervo de conhecimento assim disponibilizado.


%%% Início de seção. %%%
\section{Visualização de Grafos}
\label{sec-referencial-visualizacao}

Texto.


\subsection{Grafos e o Seu Desenho}
\label{sec-referencial-visualizacao-desenho}

Texto.

\begin{figure}[H]
    \centering
    \caption{O mesmo recorte do grafo desenhado pelo layout por forças e pelo layout hierárquico}
    \label{fig-vis-layouts}
    \captura{figuras/fig-vis-layouts.jpg}
    \notafigura{print já feito: o mesmo grafo (1º ano, ciências, ``o ano inteiro nesta matéria'') no layout de forças, na esquerda, e no hierárquico, na direita. recortei só o desenho, sem o painel e sem a legenda.}
\end{figure}


\subsection{Layout por Forças}
\label{sec-referencial-visualizacao-forcas}

Texto.


\subsection{Layout Hierárquico}
\label{sec-referencial-visualizacao-hierarquico}

Texto.

\begin{figure}[H]
    \centering
    \caption{Etapas do layout hierárquico}
    \label{fig-vis-etapas-hierarquico}
    \marcadorbrasao
    \notafigura{aqui vai um esquema simples, feito no próprio latex, com as etapas do hierárquico: separar os nós em camadas, ordenar cada camada pra cruzar menos arestas e depois calcular a posição de cada nó. podem ser três caixinhas com setas.}
\end{figure}



%%% Início de seção. %%%
\section{Acessibilidade na Web}
\label{sec-referencial-acessibilidade}

Texto.


\subsection{Diretrizes de Acessibilidade para Conteúdo Web}
\label{sec-referencial-acessibilidade-wcag}

Texto.

\begin{table}[H]
    \centering
    \caption{Critérios das WCAG atendidos pelo EduGraph}
    \label{tab:wcag}
    \renewcommand{\arraystretch}{1.3}
    \begin{tabular}{|l|l|l|}
        \hline
        \textbf{Critério} & \textbf{O que pede} & \textbf{Como o site atende} \\
        \hline
        Texto. & Texto. & Texto. \\
        \hline
    \end{tabular}
    \notafigura{aqui vai uma tabela com os critérios das wcag que o site atende: contraste, não depender só da cor, aumentar o texto, usar pelo teclado, foco visível e animação que dá pra desligar. cada linha com o número do critério, o que ele pede e onde isso aparece no site.}
\end{table}


\subsection{Deficiências na Percepção de Cores}
\label{sec-referencial-acessibilidade-cores}

Texto.

\begin{figure}[H]
    \centering
    \caption{Cores dos níveis do grafo na visão típica e nas simulações de daltonismo}
    \label{fig-acess-simulacao}
    \marcadorbrasao
    \notafigura{aqui vai uma tabela de cores, feita no latex mesmo, com a cor de cada nível do grafo (ano, área, componente, unidade, objeto e habilidade) na visão normal e simulando protanopia, deuteranopia e tritanopia. a ideia é mostrar que dá pra diferenciar as cores em todas.}
\end{figure}



%%% Início de seção. %%%
\section{Tecnologias Utilizadas}
\label{sec-referencial-tecnologias}

Esta seção apresenta as tecnologias adotadas no desenvolvimento do EduGraph e as razões de cada escolha; a Tabela~\ref{tab:tecnologias} resume-as, com as versões utilizadas, e a forma como se articulam é descrita na Seção~\ref{sec-edugraph-arquitetura}.

No \textit{front-end}, adotou-se o React, biblioteca JavaScript de código aberto voltada à construção de interfaces de usuário a partir de componentes declarativos e reutilizáveis. O modelo de componentes ajusta-se à decomposição natural da interface projetada, a barra de busca e filtros, o palco de visualização do grafo, o painel de detalhe das entidades e as sobreposições de estado, permitindo que cada unidade encapsule estrutura, aparência e comportamento próprios, o que favorece o desenvolvimento incremental e a manutenção isolada de cada parte. A renderização declarativa complementa esse modelo: a interface é descrita como função do estado da aplicação, de modo que transições como consulta em andamento, resultado disponível, resultado vazio e falha da consulta refletem-se automaticamente na tela, sem manipulação manual da árvore do documento. Pesou ainda o caráter de biblioteca e não de \textit{framework} do React, que não impõe arquitetura de aplicação, solução de roteamento ou camada de estilos, preservando as demais decisões do projeto.

Em torno do React organizam-se as demais escolhas da camada de apresentação. O Vite atua como empacotador, respondendo pelo servidor de desenvolvimento e pela geração dos artefatos de produção. A estilização emprega CSS (\textit{Cascading Style Sheets}) puro, organizado em variáveis de tema que centralizam a paleta nos modos claro e escuro de exibição, sem o recurso a \textit{frameworks} de componentes prontos --- o que assegura controle integral sobre a identidade visual construída na etapa de prototipação, descrita na Seção~\ref{sec-edugraph-prototipacao}, e sobre os recursos de acessibilidade da ferramenta. Os estilos combinam declarações em linha nos componentes, para a estrutura e para as propriedades calculadas em função do estado da aplicação, e uma folha central com as variáveis de tema e as classes responsáveis pelos estados de interação e pelas animações. A navegação entre as telas também é controlada pelo estado do React, sem biblioteca de roteamento, de modo que as únicas dependências de execução da interface são os pacotes \texttt{react} e \texttt{react-dom}. A visualização do grafo, por sua vez, é desenhada em um elemento \textit{canvas} do HTML (\textit{HyperText Markup Language}), por meio de sua interface de desenho bidimensional (Canvas 2D), com um motor próprio de simulação e de interação, sem bibliotecas de visualização de grafos, descrito na Seção~\ref{sec-edugraph-grafos}.

A integração entre as camadas reforça essas escolhas. A interface consulta a API do próprio sistema por meio de requisições HTTP (\textit{Hypertext Transfer Protocol}), com respostas em JSON (\textit{JavaScript Object Notation}), recursos nativos do navegador utilizados a partir dos componentes React, sem bibliotecas de acesso específicas. O acoplamento entre a interface e o servidor restringe-se, assim, ao contrato das consultas e das respostas, o que permitiu, nas fases iniciais do desenvolvimento, substituir o servidor por um simulador que reproduzia o contrato então previsto, viabilizando a evolução da interface antes da disponibilização dos dados.

No armazenamento, adotou-se o Neo4j, banco de dados em grafo cuja linguagem declarativa de consulta é a Cypher \cite{angles2017foundations}. O acesso ao banco se dá pela Query API, interface que executa instruções Cypher recebidas em requisições HTTP e devolve o resultado em JSON \cite{neo4j2026queryapi}. A opção por essa interface, em detrimento dos \textit{drivers} oficiais do banco, permitiu adotar em todo o sistema um único formato de resposta, o JSON tipado, no qual cada valor é acompanhado da indicação do próprio tipo, como nó, relação, caminho, inteiro, entre outros, que a API repassa sem conversão e a partir do qual a interface reconstrói o grafo de qualquer resultado de consulta.

Entre a interface e o banco, adotou-se uma API própria, escrita em Python com o \textit{framework} FastAPI e executada pelo servidor Uvicorn; a comunicação com o Neo4j usa a biblioteca httpx, e a configuração de acesso ao banco é lida de um arquivo local pela biblioteca python-dotenv. A escolha do FastAPI considerou a validação dos dados recebidos por meio de tipos, a geração automática da documentação interativa das rotas, que favorece a continuidade do projeto por outros pesquisadores, e o suporte a rotas assíncronas, adequado a serviços que respondam de forma incremental. A mesma API deverá receber os serviços previstos para o sistema --- o fórum, o repositório de ontologias e o assistente baseado em modelos de linguagem ---, cujos módulos-base já constam dela (Seção~\ref{sec-edugraph-arquitetura}).

\begin{table}[h]
    \centering
    \caption{Tecnologias utilizadas e respectivas versões}
    \label{tab:tecnologias}
    \renewcommand{\arraystretch}{1.3}
    \begin{tabular}{|>{\raggedright\arraybackslash}p{4.2cm}|>{\raggedright\arraybackslash}p{2.4cm}|>{\raggedright\arraybackslash}p{7.6cm}|}
        \hline
        \textbf{Tecnologia} & \textbf{Versão} & \textbf{Uso no projeto} \\
        \hline
        React (\texttt{react} e \texttt{react-dom}) & 18.3.1 & Construção da interface web \\
        \hline
        Vite & 5.4.21 & Servidor de desenvolvimento e geração da versão de produção da interface \\
        \hline
        Node.js & 22.22.1 & Execução das ferramentas de desenvolvimento da interface \\
        \hline
        Python & 3.14.4 & Linguagem da API \\
        \hline
        FastAPI & 0.165.0 & \textit{Framework} da API: rotas, validação dos dados e documentação \\
        \hline
        Uvicorn & 0.54.0 & Servidor que executa a API \\
        \hline
        httpx & 0.28.1 & Cliente HTTP da API para o Neo4j \\
        \hline
        python-dotenv & 1.2.4 & Leitura da configuração de acesso ao banco \\
        \hline
        Neo4j, edição Community & 2026.09.0 & Banco de dados em grafo \\
        \hline
    \end{tabular}
\end{table}

% ==============================================================================
% PG - Lorenzo Rizzi Fiorot
% Capítulo 3 - Mapeamento Sistemático da Literatura
% ==============================================================================
\chapter{Mapeamento Sistemático da Literatura}
\label{sec-msl}

Este capítulo apresenta o mapeamento sistemático da literatura conduzido neste trabalho: o protocolo (Seção~\ref{sec-msl-protocolo}), a string de busca e as fontes consultadas (Seção~\ref{sec-msl-string}), a condução e a seleção dos estudos (Seção~\ref{sec-msl-conducao}), os resultados por questão de pesquisa (Seção~\ref{sec-msl-resultados}), os trabalhos correlatos (Seção~\ref{sec-msl-correlatos}) e uma síntese (Seção~\ref{sec-msl-sintese}).


%%% Início de seção. %%%
\section{Protocolo}
\label{sec-msl-protocolo}

O objetivo geral do mapeamento sistemático conduzido neste trabalho foi levantar e caracterizar as iniciativas de ensino do Pensamento Computacional na educação básica brasileira, identificando quais áreas do conhecimento previstas na BNCC têm sido exploradas em conjunto com os pilares da computação, bem como quais abordagens pedagógicas e ferramentas têm sido empregadas nesse contexto.

% TODO(autor): o protocolo cita o período de cobertura, mas o texto não o informa. O
% README do repositório menciona publicações de 2012 a 2025; confirmar no Parsifal antes
% de escrevê-lo aqui.
Antes da execução das buscas, foi elaborado um protocolo de pesquisa formal, registrado e gerenciado com o apoio do Parsifal \cite{parsifal}, ferramenta web de código aberto de apoio à condução de revisões sistemáticas da literatura. O protocolo definiu previamente todos os parâmetros da pesquisa, incluindo as questões de pesquisa, as fontes consultadas, o período temporal de cobertura e os critérios de inclusão e exclusão dos estudos, garantindo a reprodutibilidade e a transparência do processo.

As questões de pesquisa (QP) que orientaram o mapeamento foram derivadas do objetivo geral e organizam-se em torno dos eixos de estratégias, impactos e instrumentos de avaliação do Pensamento Computacional na educação básica, conforme apresentado na Tabela~\ref{tab:questoes-pesquisa}.

\begin{table}[h]
    \centering
    \caption{Questões de pesquisa}
    \label{tab:questoes-pesquisa}
    \renewcommand{\arraystretch}{1.3}
    \begin{tabular}{|p{15cm}|}
        \hline
        \multicolumn{1}{|c|}{\textbf{Questões de Pesquisa}} \\
        \hline
        \textbf{QP1:} Quais são as estratégias pedagógicas e metodologias mais utilizadas para a introdução do Pensamento Computacional na educação básica?\\
        \textbf{QP2:} De que maneira o ensino transversal do Pensamento Computacional impacta o desenvolvimento das competências e habilidades na educação básica?\\
        \textbf{QP3:} Como as atividades de Pensamento Computacional podem ser integradas ao currículo escolar para favorecer a interdisciplinaridade?\\
        \textbf{QP4:} Quais instrumentos e critérios de avaliação são utilizados na literatura para mensurar o efetivo desenvolvimento do Pensamento Computacional em estudantes da educação básica?\\
        \textbf{QP5:} Como os materiais de apoio ao professor e os recursos de avaliação do Pensamento Computacional têm sido disponibilizados, organizados e acessados na educação básica?\\
        \hline
    \end{tabular}
\end{table}

A partir dessas questões, foram definidos os critérios de inclusão e exclusão que delimitaram o conjunto de estudos submetidos à análise. Enquanto os critérios de inclusão (CI) estabelecem as características que tornam um estudo elegível ao mapeamento --- como a aderência temática ao Pensamento Computacional, o foco na educação básica e o relato de experiências práticas ou revisões ---, os critérios de exclusão (CE) especificam as condições que motivam o descarte de trabalhos, garantindo que somente as publicações alinhadas ao escopo e aos objetivos da pesquisa fossem efetivamente consideradas, conforme apresentado na Tabela~\ref{tab:criterios-selecao}.

\begin{table}[h]
    \centering
    \caption{Critérios de seleção}
    \label{tab:criterios-selecao}
    \renewcommand{\arraystretch}{1.3}
    \begin{tabular}{|p{15cm}|}
        \hline
        \multicolumn{1}{|c|}{\textbf{Critérios de Inclusão}} \\
        \hline
        \textbf{CI1:} Estudos que abordem estratégias ou metodologias de ensino de Pensamento Computacional. \\
        \textbf{CI2:} Estudos com foco na educação básica. \\
        \textbf{CI3:} Estudos que apresentem resultados relacionados ao desenvolvimento de habilidades cognitivas, como lógica e resolução de problemas. \\
        \textbf{CI4:} Estudos que relatem experiências práticas, intervenções ou revisões sistemáticas sobre o tema. \\
        \textbf{CI5:} Estudos que apresentem, descrevam ou analisem materiais didáticos, recursos de avaliação ou instrumentos de apoio ao professor para o ensino do Pensamento Computacional na educação básica. \\
        \hline
        \multicolumn{1}{|c|}{\textbf{Critérios de Exclusão}} \\
        \hline
        \textbf{CE1:} Estudos duplicados, sem resumo ou texto completo disponível. \\
        \textbf{CE2:} Publicações que não configurem artigo científico completo, como slides, tutoriais, pôsteres, palestras ou editoriais. \\
        \textbf{CE3:} Estudos publicados apenas como resumo expandido ou \textit{short paper} com menos de quatro páginas. \\
        \textbf{CE4:} Trabalhos não redigidos em inglês ou português, ou versões mais antigas de outras publicações já selecionadas. \\
        \textbf{CE5:} Pesquisas voltadas exclusivamente ao Ensino Superior. \\
        \textbf{CE6:} Trabalhos sem aderência completa ao tema investigado. \\
        \hline
    \end{tabular}
\end{table}


%%% Início de seção. %%%
\section{String de Busca e Fontes}
\label{sec-msl-string}

A string de busca foi construída com base no \textit{framework} PICOC (\textit{Population, Intervention, Comparison, Outcome, Context}), método consolidado para a derivação sistemática de termos \cite{kitchenham2007guidelines}. A população de interesse é composta por estudantes e professores da educação básica; a intervenção corresponde ao ensino do Pensamento Computacional; o resultado esperado refere-se ao desenvolvimento de habilidades computacionais alinhadas à BNCC; e o contexto é o cenário educacional brasileiro. A partir da análise desses elementos, foram derivados os termos centrais da busca, combinados por meio dos operadores booleanos \texttt{AND} e \texttt{OR}, conforme apresentado a seguir:

\begin{verbatim}
("Pensamento Computacional" OR "Computational Thinking"
 OR "Computação Desplugada" OR "Unplugged Computing")
                            AND
("Educação Básica" OR "Basic Education" OR "Ensino Fundamental"
 OR "Elementary School" OR "BNCC")
                            AND
("Ensino" OR "Teaching" OR "Aprendizagem" OR "Learning"
 OR "Currículo" OR "Curriculum")
\end{verbatim}

A escolha priorizou duas fontes nacionais, a SBC OpenLib (SOL), biblioteca digital da Sociedade Brasileira de Computação (SBC), e a Revista Novas Tecnologias na Educação (RENOTE), dado que o foco do mapeamento é o contexto educacional brasileiro. A string foi adaptada à sintaxe específica de cada base, mantendo a equivalência semântica entre as versões.


%%% Início de seção. %%%
\section{Condução e Seleção dos Estudos}
\label{sec-msl-conducao}

% TODO(autor): faltam os números da seleção, por etapa e por fonte (estudos retornados,
% duplicados, excluídos por título e resumo, excluídos por leitura completa e
% selecionados), a exportar do Parsifal, e a lista dos estudos selecionados. Sem esses
% dados, a seção não traz números.
Texto.


%%% Início de seção. %%%
\section{Resultados}
\label{sec-msl-resultados}

Esta seção sintetiza as respostas às questões de pesquisa formuladas no protocolo, a partir da análise integrada dos estudos selecionados nas duas fontes consultadas (SOL e RENOTE). A apresentação organiza-se segundo as cinco questões de pesquisa (QP1 a QP5), reunindo, em cada uma, as evidências mais representativas do conjunto analisado.

\subsection{Estratégias Pedagógicas (QP1)}
\label{sec-msl-resultados-qp1}

Quanto às estratégias pedagógicas e metodologias (QP1), o conjunto analisado evidencia o predomínio da Computação Desplugada, sobretudo nos anos iniciais e em contextos de baixa infraestrutura tecnológica, frequentemente materializada em jogos não digitais \cite{silvamadeira2024}. Em caráter complementar, observa-se a adoção de metodologias ativas --- como Sala de Aula Invertida, gamificação e narrativas \cite{boeres2024}, o Modelo Instrucional 5E articulado ao \textit{Design Thinking} \cite{danielfilho2025} e a Robótica Educacional sob a abordagem STEAM (\textit{Science, Technology, Engineering, Arts and Mathematics}) \cite{osorio2025}. Entre as ferramentas, o Scratch destaca-se como o recurso mais recorrente \cite{kohler2025}, ao lado da articulação entre atividades plugadas e desplugadas como forma de introduzir os conceitos de modo gradual.

\subsection{Impacto no Desenvolvimento de Competências (QP2)}
\label{sec-msl-resultados-qp2}

No que se refere ao impacto do ensino transversal sobre o desenvolvimento de competências e habilidades (QP2), os estudos associam o Pensamento Computacional ao aprimoramento do raciocínio lógico e da capacidade de resolução de problemas, apontados como os ganhos mais recorrentes \cite{rege2023}. As intervenções empíricas corroboram esses efeitos: \citeonline{raiol2023} registraram que cerca de 82\% dos estudantes atingiram o nível multiestrutural de compreensão. Observa-se, contudo, um desenvolvimento desigual entre os pilares --- com a decomposição figurando como habilidade de alta dificuldade \cite{joaquim2023} --- e a dependência de fatores estruturais, com destaque para a formação docente como condição necessária à efetivação dessas competências \cite{boeres2024}.

\subsection{Integração ao Currículo (QP3)}
\label{sec-msl-resultados-qp3}

Relativamente à integração ao currículo e à interdisciplinaridade (QP3), a estratégia recorrente consiste em ancorar o Pensamento Computacional nas habilidades da BNCC, tratando-o como eixo transversal. Essa articulação manifesta-se no mapeamento explícito entre os pilares do PC e as habilidades da BNCC, organizado por disciplina \cite{osorio2025}. A programação é igualmente empregada como linguagem para a modelagem de problemas de outras áreas, mobilizando conceitos de física e matemática \cite{danielfilho2025}. Registra-se, todavia, que a transversalidade prevista nos documentos normativos ainda se concretiza, em parte, de modo disciplinarmente segmentado \cite{boeres2024}.

\subsection{Instrumentos de Avaliação (QP4)}
\label{sec-msl-resultados-qp4}

Quanto aos instrumentos e critérios de avaliação (QP4), entre as propostas mais estruturadas, destacam-se o Protocolo para Avaliação por Critérios do Pensamento Computacional (PAPC), baseado em avaliação por critérios com pesos atribuídos conforme a complexidade \cite{santos2023papc}, modelos fundamentados nas taxonomias de Bloom e SOLO (\textit{Structure of the Observed Learning Outcome}) operacionalizados por pré e pós-testes \cite{raiol2023} e o recurso a ferramentas automatizadas associadas a rubricas, como o CodeMaster \cite{joaquim2023}. Entretanto, o conjunto revela heterogeneidade de abordagens: \citeonline{kohler2025} apontam a ausência de métodos e de métricas padronizados e identificam a observação, os questionários e a análise qualitativa como os instrumentos mais frequentes.

\subsection{Disponibilização dos Materiais de Apoio (QP5)}
\label{sec-msl-resultados-qp5}

Por fim, quanto à disponibilização, organização e acesso aos materiais de apoio e aos recursos de avaliação (QP5), observa-se uma tendência crescente à adoção de repositórios e plataformas \textit{online}. As soluções mais aderentes ao escopo deste trabalho organizam repositórios web colaborativos que estruturam os recursos e os vinculam ao currículo: em um caso, projetos de robótica com micro:bit organizados em \textit{cards} e associados, por disciplina, às habilidades da BNCC e aos pilares do PC \cite{osorio2025}; em outro, um repositório de planos de aula organizado por nível de ensino e por eixo da Computação, com filtragem e disponibilização sob licença Creative Commons \cite{monteiro2024}. De modo complementar, registram-se \textit{e-books} e cursos em plataformas institucionais \cite{boeres2024, rafalski2023}.


%%% Início de seção. %%%
\section{Trabalhos Correlatos}
\label{sec-msl-correlatos}

% TODO(autor): os repositórios apontados na QP5 e outros trabalhos que o autor indicar,
% comparados ao que este trabalho entrega.
Texto.


%%% Início de seção. %%%
\section{Síntese}
\label{sec-msl-sintese}

% TODO(autor): uma síntese que percorra as cinco questões de pesquisa pode ser escrita
% junto com o fechamento do texto; abaixo fica a constatação do anteprojeto.
Constata-se, a partir dos resultados da QP5, que o acervo de materiais de apoio ao professor permanece disperso entre fontes heterogêneas, o que dificulta sua localização integrada e evidencia a pertinência de uma solução centralizada e semanticamente estruturada --- necessidade que motiva o repositório de conhecimento proposto neste trabalho.

% ==============================================================================
% PG - Lorenzo Rizzi Fiorot
% Capítulo 4 - O EduGraph
% ==============================================================================
% TODO(autor): o modelo do curso sugere que o título deste capítulo nomeie a
% contribuição do trabalho; "O EduGraph" é provisório.
\chapter{O EduGraph}
\label{sec-edugraph}

Texto.


%%% Início de seção. %%%
\section{Requisitos}
\label{sec-edugraph-requisitos}

% TODO(autor): os requisitos funcionais e não funcionais foram registrados em algum
% documento? Sem esse registro, a seção não descreve requisitos que não existam.
Texto.


%%% Início de seção. %%%
\section{Arquitetura}
\label{sec-edugraph-arquitetura}

O EduGraph organiza-se em três camadas, ilustradas na Figura~\ref{fig-arquitetura}: a interface web, executada no navegador; a API, que concentra o acesso aos dados; e o banco de grafos Neo4j, que armazena o repositório de conhecimento. Para obter os dados, o navegador comunica-se apenas com o endereço do próprio sistema: em desenvolvimento, o servidor do Vite entrega a interface e repassa à API os pedidos cujo caminho começa com \texttt{/api/}, de modo que a interface não conhece o endereço do banco nem as credenciais de acesso. Na versão atual, todos esses serviços são executados localmente; a publicação do sistema depende da definição da hospedagem.

% Diagrama feito só com caixas e setas do próprio LaTeX, sem pacotes de desenho, para não
% aumentar o tempo de compilação.
\begin{figure}[htb]
    \centering
    \caption{Arquitetura do EduGraph}
    \label{fig-arquitetura}
    \setlength{\fboxsep}{7pt}
    \renewcommand{\arraystretch}{1.2}
    \begin{tabular}{c}
        \fbox{\parbox{12cm}{\centering\textbf{Navegador}\\[2pt]
            Interface web em React: páginas, busca guiada\\
            e visualização do grafo em Canvas 2D}} \\
        $\Big\updownarrow$\ {\small HTTP: a interface e as rotas \texttt{/api/}, com respostas em JSON} \\
        \fbox{\parbox{12cm}{\centering\textbf{Servidor do Vite} (em desenvolvimento)\\[2pt]
            Entrega a interface e repassa à API os pedidos que começam com \texttt{/api/}}} \\
        $\Big\updownarrow$\ {\small HTTP: rotas da API, com respostas em JSON} \\
        \fbox{\parbox{12cm}{\centering\textbf{API em Python (FastAPI)}\\[2pt]
            Consultas da aplicação por assunto e verificação do serviço;\\
            guarda a credencial do banco e impõe o modo somente leitura;\\
            módulos-base do fórum, das ontologias e do assistente}} \\
        $\Big\updownarrow$\ {\small Query API (HTTP): Cypher em modo de leitura, resposta em JSON tipado} \\
        \fbox{\parbox{12cm}{\centering\textbf{Banco de grafos Neo4j} (edição Community)\\[2pt]
            Repositório de conhecimento: o grafo da BNCC (Seção~\ref{sec-edugraph-modelo})}} \\
    \end{tabular}
\end{figure}

O contrato entre a interface e a API baseia-se em requisições HTTP e respostas em JSON. Cada consulta da aplicação --- os números da página inicial, os anos, as áreas do conhecimento e os componentes de um ano, o recorte de um ano em um componente, as perguntas da busca guiada e as conexões de um nó --- é exposta por uma rota da API que executa uma única instrução Cypher, sempre com parâmetros, de modo que os valores informados nunca são inseridos no texto da consulta. A resposta reúne a instrução executada, os parâmetros e o resultado devolvido pelo Neo4j, sem alteração, no formato JSON tipado; com isso, a interface reconstrói os nós, as relações e os caminhos de qualquer resultado e exibe ao usuário a consulta realizada (Seção~\ref{sec-edugraph-grafos}). Os erros do banco, da API e do servidor de desenvolvimento seguem um formato único, com um código e uma mensagem, o que permite à interface tratá-los de forma uniforme.

A credencial de acesso ao banco fica apenas na configuração da API, em um arquivo local excluído do controle de versão, e nunca é enviada ao navegador. Toda consulta é encaminhada ao Neo4j em modo de leitura, independentemente do que a interface solicite, e o próprio servidor recusa instruções de escrita nesse modo, comportamento verificado durante o desenvolvimento. A medida é necessária porque o controle de acesso por papéis é recurso da edição Enterprise do Neo4j \cite{neo4j2026autorizacao}: na edição Community, adotada no projeto, não há como restringir um usuário à leitura. Por fim, a API valida os parâmetros recebidos, como os códigos de ano, de área e de componente e os limites de resultados, e limita a 12 segundos a espera por uma resposta do banco.

A API reúne, ainda, uma rota que informa se o serviço e a conexão com o banco estão disponíveis e os módulos-base de três serviços previstos --- o fórum, o repositório de ontologias e o assistente ---, cujas rotas já constam da documentação interativa gerada pelo FastAPI e respondem como ainda não implementadas (Seção~\ref{sec-edugraph-api}).


%%% Início de seção. %%%
\section{Modelo do Grafo}
\label{sec-edugraph-modelo}

% TODO(autor): dizer aqui quem definiu o modelo do grafo e quem mantém a base; a redação
% depende da divisão de responsabilidades descrita na Introdução.
O repositório de conhecimento é um grafo de propriedades armazenado no Neo4j, cujo modelo reproduz o encadeamento da BNCC descrito na Seção~\ref{sec-referencial-bncc}: ano, área do conhecimento, componente curricular, unidade temática, objeto de conhecimento e habilidade. Cada nível corresponde a nós com rótulo próprio e liga-se ao nível seguinte por uma relação nomeada, conforme a Tabela~\ref{tab:modelo-grafo}; as relações não têm propriedades. Nos três últimos níveis, o rótulo leva o sufixo do componente curricular --- \texttt{Mat}, para Matemática, e \texttt{Cien}, para Ciências ---, como em \texttt{HabilidadeMat}. Além do encadeamento, cada ano liga-se diretamente às suas habilidades pela relação \texttt{DESENVOLVE\_HABILIDADE}.

% TODO(autor): na base atual, as descrições das habilidades de Ciências coincidem com o
% texto da BNCC, e as de Matemática estão resumidas. Se continuarem assim, dizer isso aqui.
\begin{table}[h]
    \centering
    \caption{Níveis do modelo do grafo}
    \label{tab:modelo-grafo}
    \renewcommand{\arraystretch}{1.3}
    \begin{tabular}{|>{\raggedright\arraybackslash}p{4.2cm}|>{\raggedright\arraybackslash}p{5.6cm}|>{\raggedright\arraybackslash}p{4.6cm}|}
        \hline
        \textbf{Nível} & \textbf{Rótulo dos nós} & \textbf{Relação com o nível seguinte} \\
        \hline
        Ano & \texttt{Ano} & \texttt{TEM\_AREA} \\
        \hline
        Área do conhecimento & \texttt{AreaConhecimento} & \texttt{TEM\_COMPONENTE} \\
        \hline
        Componente curricular & \texttt{ComponenteCurricular} & \texttt{TEM\_UNIDADE} \\
        \hline
        Unidade temática & \texttt{UnidadeTematica} + sufixo & \texttt{TEM\_OBJETO} \\
        \hline
        Objeto de conhecimento & \texttt{ObjetoConhecimento} + sufixo & \texttt{DESENVOLVE} \\
        \hline
        Habilidade & \texttt{Habilidade} + sufixo & --- \\
        \hline
    \end{tabular}
\end{table}

Os nós de área, de componente, de unidade temática e de objeto de conhecimento não se repetem por ano: um nó é compartilhado por todos os anos em que aparece, como a unidade temática Terra e Universo, presente nos quadros do 1º e do 2º ano, à qual se ligam os objetos de conhecimento dos dois anos. O que associa cada habilidade a um ano é a relação direta \texttt{DESENVOLVE\_HABILIDADE}; as consultas da aplicação partem dela para restringir os resultados ao ano escolhido, o que impede que conteúdos de anos distintos se misturem nos subgrafos exibidos. Cada ano, por sua vez, liga-se apenas às áreas do conhecimento em que tem conteúdo.

Embora os rótulos variem conforme o componente, as consultas com que a aplicação obtém os dados percorrem o grafo pelas relações, comuns a todos os componentes, e só mencionam os rótulos \texttt{Ano}, \texttt{AreaConhecimento} e \texttt{ComponenteCurricular}. Assim, um componente carregado com a mesma estrutura passa a ser consultado sem alteração do código.

Em 26 de novembro de 2026, a base continha 75 nós e 108 relações, distribuídos conforme a Tabela~\ref{tab:cobertura-grafo}. Dos quatro nós de área do conhecimento, dois --- Ciências da Natureza e Matemática --- ligam-se a componentes curriculares; Linguagens e Ciências Humanas ainda não têm componentes carregados.

\begin{table}[h]
    \centering
    \caption{Conteúdo da base de dados em 26 de novembro de 2026}
    \label{tab:cobertura-grafo}
    \renewcommand{\arraystretch}{1.3}
    \begin{tabular}{|l|l|c|c|c|}
        \hline
        \textbf{Ano} & \textbf{Componente} & \textbf{Unidades temáticas} & \textbf{Objetos} & \textbf{Habilidades} \\
        \hline
        1º ano & Matemática & 5 & 17 & 22 \\
        \hline
        1º ano & Ciências & 3 & 3 & 6 \\
        \hline
        2º ano & Ciências & 3 & 3 & 8 \\
        \hline
        \multicolumn{2}{|l|}{Nós distintos} & 8 & 23 & 36 \\
        \hline
    \end{tabular}
\end{table}

\noindent As três unidades temáticas de Ciências são as mesmas nos dois anos, o que explica o total de oito unidades distintas. Somam-se aos nós da tabela os 2 nós de ano, os 4 de área do conhecimento e os 2 de componente curricular. Das 108 relações, 36 são do tipo \texttt{DESENVOLVE\_HABILIDADE}, uma para cada habilidade.

% TODO(autor): a base não tem nós de conceitos do Pensamento Computacional nem habilidades
% da BNCC Computação; dizer aqui como o trabalho trata isso, conforme o enquadramento
% definido na Introdução.


%%% Início de seção. %%%
\section{Prototipação}
\label{sec-edugraph-prototipacao}

Texto.


%%% Início de seção. %%%
\section{Interface e Identidade Visual}
\label{sec-edugraph-interface}

Texto.


%%% Início de seção. %%%
\section{Página Inicial}
\label{sec-edugraph-inicial}

Texto.


%%% Início de seção. %%%
\section{Página de Grafos}
\label{sec-edugraph-grafos}

Texto.


%%% Início de seção. %%%
\section{Acessibilidade}
\label{sec-edugraph-acessibilidade}

Texto.


%%% Início de seção. %%%
\section{API e Base de Dados}
\label{sec-edugraph-api}

Texto.


%%% Início de seção. %%%
\section{Telas Complementares}
\label{sec-edugraph-complementares}

% TODO(autor): as telas de conta, o Fórum e as Ontologias entram aqui como telas
% implementadas só no front-end, com conteúdo ilustrativo, ou só como trabalho futuro?
Texto.

% ==============================================================================
% PG - Lorenzo Rizzi Fiorot
% Capítulo 5 - Verificação e Resultados
% ==============================================================================
\chapter{Verificação e Resultados}
\label{sec-verificacao}

Texto.


%%% Início de seção. %%%
\section{Método de Verificação}
\label{sec-verificacao-metodo}

Texto.


%%% Início de seção. %%%
\section{Qualidade do Código}
\label{sec-verificacao-qualidade}

Texto.


%%% Início de seção. %%%
\section{O Sistema Entregue}
\label{sec-verificacao-sistema}

Texto.


%%% Início de seção. %%%
\section{Avaliação com Usuários}
\label{sec-verificacao-avaliacao}

% TODO(autor): a aplicação foi avaliada com as orientadoras, que são professoras. Há
% registro dessa avaliação (roteiro, observações)? Sem registro, ela entra como
% limitação ou trabalho futuro, sem resultados.
Texto.

% ==============================================================================
% PG - Lorenzo Rizzi Fiorot
% Capítulo 6 - Conclusão
% ==============================================================================
\chapter{Conclusão}
\label{sec-conclusoes}

Texto.


%%% Início de seção. %%%
\section{Considerações Finais}
\label{sec-conclusoes-consideracoes}

Texto.


%%% Início de seção. %%%
\section{Limitações}
\label{sec-conclusoes-limitacoes}

Texto.


%%% Início de seção. %%%
\section{Trabalhos Futuros}
\label{sec-conclusoes-trabalhosfuturos}

Texto.




%%% Páginas finais do documento: bibliografia e anexos. %%%

% Finaliza a parte no bookmark do PDF para que se inicie o bookmark na raiz e adiciona espaço de parte no sumário.
\phantompart

% Marca o início dos elementos pós-textuais.
\postextual

% Referências bibliográficas
\bibliography{bibliografia}


% Apêndices: entram quando houver conteúdo (por exemplo, o protocolo completo do
% mapeamento e a lista dos estudos selecionados). Ambiente vazio não compila bem,
% por isso fica comentado até lá.
% \begin{apendicesenv}
% \partapendices
% \include{apendices/apA-estudos}
% \end{apendicesenv}

% Fim do documento.
\end{document}
